% Point to the main file
\documentclass[../../Main.tex]{subfiles}

% ========== Start the document ==========
\begin{document}

\section{Section 7.2 - Probability Theory}

We didn't actually go over much

\subsection{Event Independence}

\begin{definition}[Event Independence]
    If $H$ and $Q$ are events, we say $H$ and $Q$ are \yellowbox{independent} if and only if

    $$P(H \cap Q) = P(H) \cdot P(Q)$$
\end{definition}

\begin{example}
    A card is selected from a standard 52 deck. Determine the probability that

    \begin{enumerate}
        \item
        {
            The card is a heart

            \begin{align*}
                P(\heartsuit) &= \frac{\text{\# hearts}}{\text{\# cards}} \\
                &= \frac{13}{52} \\
                &= \boxed{\frac{1}{4}} \\
            \end{align*}
        }
        \item 
        {
            The card is a queen

            \begin{align*}
                P(\text{Queen}) &= \frac{\text{\# queens}}{\text{\# cards}} \\
                &= \frac{4}{52} \\
                &= \boxed{\frac{1}{13}} \\
            \end{align*}
        }
        \item 
        {
            The card is a heart and a queen

            \begin{align*}
                P(\heartsuit \text{ and queen}) &= \frac{\text{\# heart and queen}}{\text{\# cards}} \\
                &= \frac{1}{52} 
            \end{align*}

            Notice that 

            \begin{align*}
                P(\heartsuit \text{ and queen}) &= P(\heartsuit) \cdot P(\text{queen}) \\
            \end{align*}

            This means that $H$ and $Q$ are \textbf{independent}
        }
    \end{enumerate}
\end{example}

I would put the examples leading up to the derivation of this next thing but I'm too lazy

\subsection{Binomial Probability}

\begin{definition}[Binomial Probability]
    Given a limited event (ie 2 outcomes) with probability $p$, Then 

    $$P(k \text{ successes in } n \text{ trials}) = {n \choose k} \; p^k \cdot (1 - p)^{n - k}$$
\end{definition}

\begin{example}
    Let's do some examples

    A biased coin has probability of 0.51 for landing heads when flipped. Determine the probability of the coin landing on heads exactly 6 times when the coin is flipped 10 times

    \begin{align*}
        p &= 0.51 \\
        k &= 6 \tag{successes} \\
        n &= 10 \\
        \\
        P(6 \text{ successes in } 10 \text{ trials}) &= {10 \choose 6} \; (0.51)^6 \cdot (1 - 0.51)^{10 - 6} \\
        &\approx \boxed{0.213}
    \end{align*}
\end{example}

\begin{example}
    Let's do another

    Determine the probability of rolling a 12 \textbf{at least once} when a pair of dice is rolled 20 times.

    Since we have "at least" let's do the opposite

    \begin{align*}
        P(\text{never rolling 12 out of 20}) &= {20 \choose 0} \; \left(\frac{1}{36}\right)^0 \cdot \left(\frac{35}{36}\right)^{20} \\
        &\approx 0.569 \ldots \\
    \end{align*}

    \begin{align*}
        P(\text{rolling 12 at least once out of 20}) &= 1 - P(\text{never rolling 12 out of 20}) \\
        &\approx 1 - 0.569 \\
        &\approx \boxed{0.431} \\
    \end{align*}

\end{example}

% ========== End the document ==========
\end{document}